\documentclass[11pt,a4paper]{article}
\usepackage[margin=25mm]{geometry}
\usepackage{amsmath,amssymb,amsthm,hyperref,booktabs}
\setlength{\emergencystretch}{2em}
\newtheorem{theorem}{Theorem}
\newtheorem{lemma}[theorem]{Lemma}
\newtheorem{proposition}[theorem]{Proposition}
\newcommand{\F}{\mathbb F}
\newcommand{\gen}[1]{\langle #1\rangle}
\title{Cycle counts in biaffine and Suzuki incidence graphs}
\author{Human Compression project: L025 recovery}
\date{}
\begin{document}\maketitle
\begin{abstract}
We recover the cycle counts and small graph constructions in the fixed
2021 preprint of Araujo-Pardo and Leemans. The Suzuki calculation reuses
the project's existing proof and its partial formal certificate. Two
different graphs occur in the preprint's coordinate and adjacency-list
descriptions. We distinguish them and prove the stated extremality of
both. A separate shortest-odd-cycle argument supplies the missing
unrestricted biaffine extremality proof. Its degree-four endpoint retains
an explicitly named finite classification. The complete source contract
distinguishes the recovered results, corrected statements, borrowed
mathematics and previously recorded conjecture answer.
\end{abstract}

\section{Cycle counting and the comparison class}
An $\operatorname{egr}(v,k,g,\lambda)$ is a finite simple $k$-regular graph
with $v$ vertices, girth $g$, and exactly $\lambda$ girth cycles through
each edge. Cycles are unoriented. A cycle through a fixed directed edge
has exactly one traversal starting with that directed edge. This convention
will be used throughout; it introduces no factor of two.

\begin{lemma}\label{bounds}
Let $n_0(k,2s+1)=1+k\sum_{j=0}^{s-1}(k-1)^j$ and
$n_0(k,2s)=2\sum_{j=0}^{s-1}(k-1)^j$ be the Moore tree orders.
For odd girth $g=2s+1\geq3$,
\[
 v\geq n_0(k,g)+(k-1)^s-\lambda.
\]
For even girth $g=2s\geq4$, put $D=(k-1)^s-\lambda$. Then
\[
 v\geq n_0(k,g)+\left\lceil\frac{2D}{k}\right\rceil.       \tag{1}
\]
If the even-girth graph is bipartite, the stronger bound is
\[
 v\geq n_0(k,g)+2\left\lceil\frac{D}{k}\right\rceil.        \tag{2}
\]
In particular the first formula at girth five is
$v\geq1+k^2+(k-1)^2-\lambda$; at girth six, the common
Moore term in (1)--(2) is $2(k^2-k+1)$ and $D=(k-1)^3-\lambda$.
\end{lemma}
\begin{proof}
At odd girth $2s+1$, the ball of radius $s$ about a vertex has
$n_0(k,2s+1)$ vertices: a collision in the breadth-first tree before
this radius would give a shorter cycle. Its only possible extra edges
join two vertices of depth $s$ whose first branches are distinct.
Each such edge closes one girth cycle through the root, and every
such cycle is obtained once. There are $k\lambda/2$ of them, since
a cycle through the root uses two incident edges. The $k(k-1)^{s-1}$
leaves have $k(k-1)^s$ free incidences after removing their parent
edges. Internal edges consume $k\lambda$, leaving
$k((k-1)^s-\lambda)$ edges to vertices outside the ball. Each outside
vertex receives at most $k$, proving the odd-girth bound.

At even girth $2s$ start at an edge $uv$. Its two trees of depth
$s-1$ are disjoint and have $n_0(k,2s)$ vertices. A collision,
an edge within one tree other than its tree edges, or an additional
edge touching a nonleaf would give a cycle shorter than $2s$.
Thus the only additional internal edges can join leaves of the two
trees. Each closes exactly one girth cycle through $uv$, so there
are $\lambda$ such edges. On each side the leaves have $(k-1)^s$
free incidences after removing their parent edges. Exactly $\lambda$
are used by cross edges, and the remaining $D$ go outside the trees.
Thus $2D$ edges go to at least $\lceil2D/k\rceil$ outside vertices,
proving (1).

In the bipartite case the two leaf sets lie in opposite classes.
Their outside neighbors therefore lie in disjoint classes. Each
class contains at least $\lceil D/k\rceil$ additional vertices,
giving (2).
\end{proof}
These are the bounds of Drglin, Filipovski, Jajcay and Raiman
used by the source, with their counting argument supplied here.
The fixed preprint prints exponent $(g-1)/2$ in the even-girth
numerators and in the hypothesis of Theorem 4.2. The leaf count
requires $g/2$, as in (1)--(2); it is this value that gives the
source's subsequent numerical substitution. The two comparison
classes in (1) and (2) must also be kept distinct.

\section{Biaffine graphs and added matchings}
Use two copies of $\F_q^2$, whose elements are points $(x,y)$ and lines
$[m,b]$. A point and line are adjacent when $y=mx+b$. This is the
type-one biaffine graph $B_q$. Its order, degree, girth and diameter
are $2q^2,q,6,4$ for $q\geq3$. These incidence facts and their full
proofs are retained from \emph{Cages and biaffine planes}, the common
paper of the L013 recovery. In particular two points of different first
coordinates have a unique common line, while two distinct points of the
same first coordinate have none. The corresponding dual statements hold
for lines. That whole common paper is a dependency, not a free citation
to a result whose hypotheses remain unchecked.

\begin{theorem}\label{biaffine}
For every prime power $q\geq3$, the graph $B_q$ is
\[
 \operatorname{egr}\bigl(2q^2,q,6,(q-1)^2(q-2)\bigr).
\]
It has least order among all graphs with these parameters, subject at
$q=4$ to the retained finite completeness result stated below.
\end{theorem}
\begin{proof}
Fix an edge $P_1L_1$. Choose a second point $P_2$ on $L_1$ different
from $P_1$, and a second line $L_2$ through $P_2$ different from $L_1$.
There are $(q-1)^2$ choices. On $L_2$, exclude the point $P_2$ and the
unique point with the same first coordinate as $P_1$. These exclusions
are distinct. Each of the remaining $q-2$ points $P_3$ has a unique
common line $L_3$ with $P_1$. The lines $L_1,L_2,L_3$ are distinct:
an equality would give two intersections of distinct lines, or put
$P_3$ in the excluded position. The resulting hexagon
$P_1L_1P_2L_2P_3L_3$ is simple. Conversely every hexagon through the
directed edge $P_1L_1$ is obtained exactly once. This proves the count.

Here $D=(q-1)^2$. Formula (2) gives
\[
 2(q^2-q+1)+2\left\lceil\frac{(q-1)^2}{q}\right\rceil=2q^2.
\]
This is attained by $B_q$, proving bipartite extremality. Formula (1)
gives $v\geq2q^2-1$. For odd $q$, the handshaking identity $qv=2|E|$
makes $v$ even, so $v\geq2q^2$ as well.

For even $q\geq8$, use the complete companion note
\emph{A shortest-odd-cycle reduction for biaffine extremality}.
It proves that a nonbipartite $k$-regular graph of girth at least six
has at least $\lceil7(k^2-k+1)/3\rceil$ vertices. A possible
competitor of order $2q^2-1$ would be nonbipartite by (2), but
\[
7(q^2-q+1)-3(2q^2-1)=(q-5)(q-2)>0.
\]
This excludes it. At $q=4$, the retained exhaustive result of
Goedgebeur and Jooken, Table~2 of arXiv:2401.08271v2, gives
$n(4,6,18)=32$. Its lower bound excludes the remaining order 31;
we neither rerun its census nor claim its uniqueness conclusion here.
These are all even prime powers at least four, completing the proof.
\end{proof}
At $q=2$, the common coordinate description gives $B_2=C_8$.
It too is extremal, now with parameters $(v,k,g,\lambda)=(8,2,8,1)$:
any graph containing an eight-cycle has at least eight vertices.
Thus the source's extremality sentence can include $q=2$, while its
displayed girth-six family must retain $q\geq3$.

\begin{lemma}\label{matchingcount}
Let $q$ be an even prime power, and let $a,b\in\F_q^*$ be distinct.
Add to $B_q$ the matchings
\[
 (x,y)\sim(x,y+a),\qquad [m,c]\sim[m,c+b].
 \tag{3}
\]
The resulting graph has girth five. Each added edge belongs to
$q(q-1)$ pentagons, and each original incidence edge belongs to
$4(q-1)$ pentagons.
\end{lemma}
\begin{proof}
The two matchings together form a matching on the entire vertex set.
There are no triangles or quadrilaterals by the matching instance of
the common biaffine theorem: a mixed quadrilateral would require $a=b$.
For completeness, a triangle with an added edge would require a common
line of parallel points or a common point of parallel lines. A
quadrilateral with two added edges would have its two incidence
equations differ by both $a$ and $b$, again impossible. A cycle with
no added edge lies in $B_q$ and has length at least six.

A pentagon has an odd number of added edges because its other edges
switch the two types. A matching contributes at most two edges to a
pentagon, so it contributes exactly one. For an added point edge $PQ$,
choose one of the $q$ lines through $Q$, then one of its $q-1$ points
with first coordinate different from that of $P$. The unique line
joining that point to $P$ completes the pentagon. All its vertices are
distinct by the incidence properties already stated, and the choices
are recovered from the pentagon. This gives $q(q-1)$; the dual proof
handles an added line edge.

Fix an incidence edge $PL$. There are four possible positions for the
unique added edge among the other four edges of a pentagon traversed
from $P$ to $L$. If it is adjacent to $P$, its other endpoint is fixed
by (3). Choose the other point on $L$, excluding $P$: there are $q-1$
choices, and the unique joining line completes the cycle. The case
adjacent to $L$ is dual. If the added edge is opposite and of line
type, choose the other line through $P$, excluding $L$. Its matching
partner is then fixed, and meets $L$ in a unique point because their
slopes differ. There are again $q-1$ choices. The opposite point case
is dual. The four cases are disjoint and exhaust the possibilities.
Their total is $4(q-1)$, which is positive, proving girth five.
\end{proof}

For $q=4$, take $\F_4=\F_2[\alpha]/(\alpha^2+\alpha+1)$,
$a=1$ and $b=\alpha$. The two numbers in Lemma~\ref{matchingcount}
are both twelve. This gives the coordinate graph $H$ of source
Theorem 2.2. Its vertex transitivity can also be proved explicitly.
For $s,t,c\in\F_4$, the maps
\[
 (x,y)\longmapsto(x+s,y+cx+t),\qquad
 [m,b]\longmapsto[m+c,b+t-(m+c)s]
\]
preserve incidence and both matchings. They act transitively on each
type: points can be moved using $s,t$, and lines using $c,t$.
The map interchanging the types is
\[
 (x,y)\longmapsto[x^2,\alpha y^2],\qquad
 [m,b]\longmapsto(\alpha m^2,\alpha b^2).
\]
Indeed $y=mx+b$ is equivalent to
$\alpha b^2=x^2(\alpha m^2)+\alpha y^2$. An added point
difference of one becomes a line difference of $\alpha$, and an
added line difference of $\alpha$ becomes a point difference of
$\alpha\alpha^2=1$. Thus it is an automorphism and proves vertex
transitivity. This supplies the field squaring missing from the
preprint's displayed type-interchanging map; without it, an added
line edge would map to point difference $\alpha^2$, not one.

\section{The two graphs on thirty-two vertices}
Let $G$ denote instead the literal thirty-two-row adjacency list in
source Section 2.3. Its rows are retained as exact finite data. The
following finite conclusions have direct exhaustive checks, described
below. The literal data and complete verification programs accompany
this paper; independent acceptance is recorded separately.
\[
\begin{array}{c|rrrrr}
 &v&k&g&\lambda&|\operatorname{Aut}|\\\hline
 H&32&5&5&12&1920\\
 G&32&5&5&12&48
\end{array}
\]
In particular they are not isomorphic. The fixed preprint identifies
them, whereas the abstract of the published 2022 edition explicitly
describes two graphs. We retain the fixed preprint as the intake
edition and record this distinction, without attributing the same
identification claim to the published edition.

\begin{proposition}\label{small}
Both $H$ and $G$ have least order among $\operatorname{egr}(v,5,5,12)$
graphs, subject to the explicitly retained classification of the four
$(5,5)$-cages and the finite checks below.
\end{proposition}
\begin{proof}
Lemma~\ref{bounds} gives $v\geq30$. The handshaking identity excludes
odd $v$. The retained complete classification states that there are
exactly four $5$-regular graphs of girth five on thirty vertices.
The literal graphs, in Brouwer's order, have these edge distributions:
\[
\begin{array}{c|rrrr}
 &\lambda_e=12&\lambda_e=13&\lambda_e=14&\lambda_e=16\\\hline
 1&45&0&30&0\\
 2&50&20&5&0\\
 3&60&0&15&0\\
 4&24&48&0&3
\end{array}
\]
Each row sums to seventy-five edges, and each contains an edge with
count different from twelve. None has the required constant count.
Thus $v\geq32$, attained by both displayed graphs.
\end{proof}

The finite counting procedure is elementary. For each root, enumerate
simple paths on five vertices whose other vertices exceed that root.
Accept a path when its last vertex is adjacent to its root, and its
second vertex is smaller than its last. The root is then the minimum
vertex of the cycle, and the last comparison chooses exactly one of its
two directions. Thus every pentagon is accepted once. Increment the
counter of each of its five edges. Symmetry of every adjacency list,
degree, girth and connectedness are checked independently from these
counters. The four source graphs were retrieved from the exact data
links of Brouwer's cage table; their mathematical adjacency, rather
than descriptive website metadata, is the computational input.

To establish a full automorphism order, the accepted L013 paired
partition procedure is reused in its entirety. It simultaneously
refines a source and target partition by the number of neighbors in
each cell. A disagreement excludes a partial map. Otherwise choose
one source vertex in a nonsingleton cell and try every target in the
corresponding cell; refine again. At a discrete partition, check the
entire adjacency relation literally. Every automorphism follows one
branch, every accepted leaf is an automorphism, and different leaves
give different maps. The search therefore proves the full order,
rather than merely the order of a discovered subgroup.

The group of $G$ has the source's structure $E_{16}:C_3$.
The saved permutation generators preserve adjacency literally;
their closure has forty-eight elements, hence is the full group
by the preceding independent exhaustive count. The identity and
all elements of order two form a closed abelian subgroup $E$ of
order sixteen. A recorded element $c$ of order three normalizes
$E$, and $E,Ec,Ec^2$ exhaust the group. These exact product checks
prove the asserted semidirect product. The recorded action has
four fixed elements on $E$, so the action is specified as well
as the abstract factors. Discovery of the generators is not an
extra assumption in their literal verification.

The four-cage completeness theorem is retained external mathematics.
The local data and cycle checks do not reprove that classification.
Their program lengths and the historical source programs remain
separate from a manuscript-only comparison. No factor is assigned
until the complete recovery has been reviewed and measured.

\section{The Suzuki graph as a coset walk}
We now reuse the project's existing Suzuki pentagon proof. Let
$q=2^{2e+1}$, $e\geq1$, and $\theta=2^{e+1}$. Thus applying
$x\mapsto x^\theta$ twice gives $x\mapsto x^2$ on $\F_q$.
Use the standard faithful Suzuki representation, with
\[
 z(a)=\begin{pmatrix}
 1&0&0&0\\0&1&0&0\\a&0&1&0\\a^\theta&a&0&1
 \end{pmatrix},\quad
 t(s)=\operatorname{diag}(s^{1+\theta},s,s^{-1},s^{-1-\theta}),
 \quad w=\operatorname{antidiag}(1,1,1,1).
\]
Set $Z=\{z(a):a\in\F_q\}$, $T=\{t(s):s\in\F_q^*\}$ and
$H=ZT$. The standard order formula
$|\operatorname{Sz}(q)|=q^2(q^2+1)(q-1)$, its doubly transitive
ovoid action, and the identification of its point and pair stabilizers
are retained structural inputs. They are not consequences of a
finite list of small-field checks.

The displayed matrices give $z(a)z(b)=z(a+b)$,
$H\cap wHw=T$, and $w$ normalizes $T$ by inversion.
The source's ordered nucleus and circle-point pair can be moved to
the standard ordered ovoid pair by double transitivity. Its elementary
abelian subgroup of order $q$ is $Z$: the involutions of the point
stabilizer's Sylow $2$-subgroup are exactly its center. The common
two-point stabilizer is then $T$, giving $H=ZT$. Its unordered pair
stabilizer is $K=\gen{T,w}$. These are the source's exact two
parabolics, up to conjugacy. Inverting right cosets converts its
incidence graph construction to the following left-coset convention.

Vertices are $gH$ and neighbors of $gH$ are
\[
                    gz(a)wH\quad(a\in\F_q).              \tag{4}
\]
Indeed $K/T$ has two elements, so an incident $K$-coset gives an
edge with these endpoints. The $q$ neighbors in (4) are distinct:
equality would put $z(a+b)$ in $H\cap wHw=T$, forcing $a=b$.
Since $w\notin H$, there are no loops. Since $w^2=1$, adjacency
is symmetric. The group is vertex transitive and $H$ is transitive
on its neighbors, so it is arc transitive. The number of vertices
is $[\operatorname{Sz}(q):H]=q(q^2+1)$.

\begin{theorem}\label{suzuki}
The source's Suzuki graph is
\[
                  \operatorname{egr}(q(q^2+1),q,5,q-1).
\]
\end{theorem}
\begin{proof}
Fix the arc $H,wH$. Successive neighbor coordinates are unique by
(4). The next coordinate zero is precisely the backtracking choice,
because $wz(0)w=1$. Thus nonbacktracking walks of length five
starting at this arc correspond exactly to $a,b,c,d\ne0$.
Their final coset is $W(a,b,c,d)H$, where
\[
       W=wz(a)wz(b)wz(c)wz(d)w.
\]
The walk closes if and only if $W\in H$.

Put $A=a^\theta,B=b^\theta,C=c^\theta,D=d^\theta$.
Direct multiplication in characteristic two gives
\[
 W_{32}=bc+1,\qquad W_{12}=ABc+abc+a+c,
 \qquad W_{43}=BCd+bcd+b+d.                              \tag{5}
\]
All three entries vanish in every matrix of $H$. The first forces
$bc=1$. Substitution in the second gives $(AB+1)c=0$, so $AB=1$.
Injectivity of Frobenius gives $ab=1$, and hence $c=a$.
Now $BC=1$, so the third expression is $b+d$, forcing $d=b$.
Closure therefore requires
\[
               (a,b,c,d)=(a,a^{-1},a,a^{-1}).             \tag{6}
\]
Conversely multiplication of all sixteen entries gives
\[
 W(a,a^{-1},a,a^{-1})=
 \begin{pmatrix}
 A&0&0&0\\0&a^2/A&0&0\\A/a&0&A/a^2&0\\1&a/A&0&1/A
 \end{pmatrix}
              =z(a^{-1})t(a^{2-\theta})\in H.            \tag{7}
\]
For example $(a^{2-\theta})^{1+\theta}=a^\theta$, since
$\theta^2$ acts as squaring. This also verifies the less immediate
diagonal exponent in (7); the other entries follow directly.
Thus exactly $q-1$ nonbacktracking closing walks have been found.

There are no triangles: the $(3,2)$ entry of
$wz(a)wz(b)w$ is one, whereas that entry vanishes in $H$.
There are no quadrilaterals: the $(1,2)$ entry of
$wz(a)wz(b)wz(c)w$ is $a^\theta b$, which cannot vanish
for a nonbacktracking walk. A closed nonbacktracking walk of length
five in a simple graph without triangles or quadrilaterals is a
simple pentagon. Each unoriented pentagon through the fixed arc
has exactly one traversal starting with that arc. Equations
(6)--(7) therefore give $q-1$ pentagons per edge. Arc transitivity
extends the count to every edge and proves the theorem.
\end{proof}

Connectivity also follows from the retained maximality of $K=N(T)$
in $\operatorname{Sz}(q)$: the group $\gen{H,w}$ properly contains
$K$ and hence is the whole group. This structural assertion is not
needed for the cycle-count calculation itself.

\section{Attribution and the recovery contract}
The biaffine construction, both thirty-two-vertex examples and the
Suzuki coset family are the source authors' constructions. The
biaffine hexagon argument is their counting argument. The matching
count above makes explicit a common argument for all even fields;
the needed girth mechanism is retained from the L013 common theory.
The Suzuki proof is the project's earlier coset-walk calculation,
reorganized here without claiming another conjecture answer.

The existing file \texttt{SuzukiPentagons.lean} proves the matrix
closure equivalence under the characteristic-two and Frobenius
hypotheses, and the short-walk entry obstructions. Its saved report
uses only \texttt{propext}, \texttt{Classical.choice} and
\texttt{Quot.sound}. It does not formalize the finite-field
instantiation, identification of the source parabolics, or bijection
between closing walks and graph cycles. Those interfaces are written
above and were addressed in the earlier independent audit. The
symbolic program checks (5), all entries of (7), and the short-walk
obstructions over the appropriate characteristic-two Laurent
polynomial ring. The old $q=8$ permutation-coset computation is
supplementary validation. None is presented as a new Lean result.

The literal preprint's Theorem 4.3 applies a bipartite bound to
claim unrestricted extremality. Theorem~\ref{biaffine} and the
complete odd-cycle companion supply the unrestricted conclusion.
The degree-four external census is an explicit additional input;
the elementary bound alone reaches only 31 at that degree.
The result is recovered, while the gap in the source proof is recorded.
No counterexample to its correctly ranged extremality assertion is claimed.

The source's final Suzuki excess bound now follows unconditionally
from Theorem~\ref{suzuki} and Lemma~\ref{bounds}:
\[
n(q,5,q-1)\geq2q^2-3q+3,\qquad
q(q^2+1)-n(q,5,q-1)\leq q^3-2q^2+4q-3.
\]
Thus its numerical $q=8$ case is included in the same calculation.

The source Appendix offers a general test for edge-girth regularity.
Its final multiplicity set ranges only over edges found on girth cycles,
so zero-count edges can be omitted. For example, the disjoint union of
$K_4$ and the cube is cubic of girth three: its six $K_4$ edges each
lie on two triangles, while its twelve cube edges lie on none. The
displayed routine would accept it. Our routine initializes a counter
for every edge and includes zero values in the final comparison.
It enumerates girth cycles by the same least-vertex and orientation
rule described above, with the computed girth replacing five.
This provides the intended test for any finite regular graph having
finite girth, and explicitly rejects a graph without cycles. The
counterexample and all six source graphs are retained finite controls.

The source's recorded Conjectures 3.2, 4.5 and 4.6 retain their
separate question records. Theorem~\ref{suzuki} reuses the existing
answer to 3.2. This recovery makes no new claim to settle the
existence families in 4.5, and does not reopen the already recorded
later disproof of 4.6. The attached inventory records its finite calculations, retained external
inputs, complete program dependencies and partial formal certificate.
No numerical compression factor or new conjecture credit is assigned here.

\begin{thebibliography}{9}
\bibitem{GJ} J. Goedgebeur and J. Jooken, \emph{Exhaustive generation
of edge-girth-regular graphs}, arXiv:2401.08271v2, 25 June 2024, Table~2.
\bibitem{AL} G. Araujo-Pardo and D. Leemans,
\emph{Edge-girth-regular graphs arising from biaffine planes and Suzuki
groups}, arXiv:2108.06636v1, 15 August 2021; published in
Discrete Mathematics 345 (2022), 112991. The fixed intake is the
seventeen-page preprint.
\bibitem{B} A. E. Brouwer, \emph{Cages}, graph data and table,
\url{https://aeb.win.tue.nl/graphs/cages/cages.html}; four
\texttt{cage-k5.g5-0i.dr} files, $i=0,1,2,3$, retrieved 12 October 2026.
\bibitem{BB} J. N. Bray and H. B\"a\"arnhielm,
\emph{A new method for recognising Suzuki groups},
arXiv:1706.05697, standard matrix representation.
\bibitem{DFJR} A. Drglin, S. Filipovski, R. Jajcay and T. Raiman,
\emph{Extremal edge-girth-regular graphs}, Graphs and Combinatorics,
\url{https://doi.org/10.1007/s00373-021-02368-9}.
\end{thebibliography}
\end{document}
